\noindent\textbf{Resumen.} Esta parte trata de la plataforma en AWS donde vive el histórico de SECOP~II y de cómo se mantiene al día. Solo usa servicios administrados: Lambda, Step Functions, S3, Glue con Apache Iceberg, Athena y EventBridge. Cada servicio tiene los permisos justos, el gasto se frena solo al llegar a 15~USD y cada actualización semanal cuesta unos 0,07~USD. Glue costó 0,83~USD en todo el proyecto, y los resultados se publican en una web que se alimenta de S3.

% -----------------------------------------------------------------------------
\subsection{Introducción}\label{cloud:sec:problema}
Una veeduría o un periodista que quiera saber, por ejemplo, qué parte se contrata sin competencia en su departamento no tiene dónde consultarlo. Responder esa pregunta exige el histórico completo y actualizado. La Sección~\ref{sec:intro} presenta el problema, los objetivos y el ODS~16.

Esta materia se encargó de montar en AWS, sin servidores propios, una plataforma que cargue el histórico, lo actualice cada semana sin intervención, publique los resultados en una web y no se salga del presupuesto.

% -----------------------------------------------------------------------------
\subsection{Arquitectura en la nube}\label{cloud:sec:arquitectura}
Los servicios y sus conexiones están en la Fig.~\ref{cloud:fig:arquitectura}. Los datos pasan por las tres capas del patrón Medallion~\cite{databricks_medallion}, guardadas en un solo bucket de S3 (Tabla~\ref{cloud:tab:capas}).

\begin{figure*}[!t]
\centering
\resizebox{\textwidth}{!}{% Arquitectura del sistema en AWS. La usan el informe completo (partes/parte_cloud.tex) y el póster
% (arquitectura_poster.tex). Necesita tikz con arrows.meta, positioning, fit, backgrounds y calc, los
% colores bronce, plata y oro, y los íconos de ../informe_cloud/iconos/.
\begin{tikzpicture}[
  font=\scriptsize,
  >={Stealth[length=2.2mm]},
  ico/.style={inner sep=0pt},
  lbl/.style={below=2pt, align=center, font=\scriptsize, text width=2.0cm},
  flujo/.style={->, line width=0.8pt, draw=black!85},
  ctrl/.style={->, dashed, line width=0.6pt, draw=black!55},
  capa/.style={draw=#1, line width=1pt, fill=#1!10, rounded corners=3pt, inner sep=3pt},
]
\def\W{9mm}
\def\dx{2.5}
% Fila de datos: de la fuente pública a la capa oro
\node[ico] (api)    at (-0.15*\dx,0)  {\includegraphics[width=\W]{iconos/internet.pdf}};
\node[ico] (lambda) at (1*\dx,0)  {\includegraphics[width=\W]{iconos/lambda.pdf}};
\node[ico] (bronce) at (2*\dx,0)  {\includegraphics[width=\W]{iconos/s3.pdf}};
\node[ico] (glue1)  at (3*\dx,0)  {\includegraphics[width=\W]{iconos/glue.pdf}};
\node[ico] (plata)  at (4*\dx,0)  {\includegraphics[width=\W]{iconos/s3.pdf}};
\node[ico] (glue2)  at (5*\dx,0)  {\includegraphics[width=\W]{iconos/glue.pdf}};
\node[ico] (oro)    at (6*\dx,0)  {\includegraphics[width=\W]{iconos/s3.pdf}};

\node[lbl] (l-api)    at (api.south)    {Fuente pública\\SECOP~II};
\node[lbl] (l-lambda) at (lambda.south) {Lambda\\de descarga};
\node[lbl] (l-bronce) at (bronce.south) {\textbf{Bronce}\\datos crudos};
\node[lbl] (l-glue1)  at (glue1.south)  {Glue:\\limpieza};
\node[lbl] (l-plata)  at (plata.south)  {\textbf{Plata}\\datos limpios};
\node[lbl] (l-glue2)  at (glue2.south)  {Glue:\\indicadores};
\node[lbl] (l-oro)    at (oro.south)    {\textbf{Oro}\\seis resúmenes};

% Control
\node[ico] (eb)  at (1*\dx,2.6) {\includegraphics[width=\W]{iconos/eventbridge.pdf}};
\node[ico] (sfn) at (3*\dx,2.6) {\includegraphics[width=\W]{iconos/stepfunctions.pdf}};
\node[lbl] (l-eb)  at (eb.south)  {Programador\\cada lunes, 6:00};
\node[lbl, text width=3.2cm] (l-sfn) at (sfn.south) {Step Functions\\coordina los pasos};

% Consulta y seguridad
\node[ico] (catalog) at (6*\dx,-2.75) {\includegraphics[width=\W]{iconos/glue_catalog.pdf}};
\node[ico] (athena)  at (5*\dx,-2.75) {\includegraphics[width=\W]{iconos/athena.pdf}};
\node[lbl, text width=1.9cm] (l-catalog) at (catalog.south) {Catálogo\\de tablas};
\node[lbl] (l-athena)  at (athena.south)  {Athena\\consultas};

\node[ico] (iam) at (1*\dx,-2.75) {\includegraphics[width=\W]{iconos/iam.pdf}};
\node[ico] (cw)  at (2*\dx,-2.75) {\includegraphics[width=\W]{iconos/cloudwatch.pdf}};
\node[ico] (bud) at (3*\dx,-2.75) {\includegraphics[width=\W]{iconos/budgets.pdf}};
\node[lbl, text width=1.9cm] (l-iam) at (iam.south) {IAM\\permisos justos};
\node[lbl, text width=1.9cm] (l-cw)  at (cw.south)  {CloudWatch\\registros};
\node[lbl, text width=1.9cm] (l-bud) at (bud.south) {Budgets\\tope de gasto};

% Publicación: de oro a la web pública
\node[ico] (exporta) at (5*\dx,-5.35)   {\includegraphics[width=\W]{iconos/lambda.pdf}};
\node[ico] (s3web)   at (3.5*\dx,-5.35) {\includegraphics[width=\W]{iconos/s3.pdf}};
\node[ico] (cdn)     at (2*\dx,-5.35)   {\includegraphics[width=\W]{iconos/cloudfront.png}};
\node[ico] (usr)     at (-0.15*\dx,-5.35) {\includegraphics[width=\W]{iconos/users.pdf}};
\node[lbl] (l-exporta) at (exporta.south) {Lambda:\\exporta a JSON};
\node[lbl] (l-s3web)   at (s3web.south)   {\textbf{S3 web}\\sitio y datos};
\node[lbl] (l-cdn)     at (cdn.south)     {CloudFront\\HTTPS y caché};
\node[lbl] (l-usr)     at (usr.south)     {Ciudadanía\\y veedurías};

\begin{scope}[on background layer]
  \node[capa=bronce, fit=(bronce)(l-bronce)] {};
  \node[capa=plata,  fit=(plata)(l-plata)] {};
  \node[capa=oro,    fit=(oro)(l-oro)] {};
\end{scope}

\draw[flujo] (api) -- (lambda);
\draw[flujo] (lambda) -- (bronce);
\draw[flujo] (bronce) -- (glue1);
\draw[flujo] (glue1) -- (plata);
\draw[flujo] (plata) -- (glue2);
\draw[flujo] (glue2) -- (oro);
\draw[flujo] (l-oro.south) -- (catalog.north);
\draw[flujo] (catalog) -- (athena);
\draw[flujo] (l-athena.south) -- (exporta.north);
\draw[flujo] (exporta) -- (s3web);
\draw[flujo] (s3web) -- (cdn);
\draw[flujo] (cdn) -- (usr);

\draw[ctrl] (eb) -- node[above, font=\scriptsize]{inicia} (sfn);
\coordinate (bus) at ($(l-sfn.south)+(0,-0.25)$);
\coordinate (derecha) at ($(oro.east)+(1.05,0)$);
\draw[dashed, line width=0.6pt, draw=black!55] (l-sfn.south) -- (bus);
\draw[dashed, line width=0.6pt, draw=black!55] (lambda.north |- bus) -- (derecha |- bus) -- (derecha |- exporta);
\draw[ctrl] (derecha |- exporta) -- (exporta.east);
\draw[ctrl] (lambda.north |- bus) -- (lambda.north);
\draw[ctrl] (glue1.north |- bus) -- (glue1.north);
\draw[ctrl] (glue2.north |- bus) -- (glue2.north);

\coordinate (arriba) at ($(eb.north)+(0,0.55)$);
\coordinate (izq) at ($(lambda.west)+(-0.6,0)$);
\begin{scope}[on background layer]
  \node[draw=black!80, line width=0.8pt, inner sep=8pt, fit=(eb)(l-sfn)(lambda)(oro)(l-oro)(l-catalog)(l-iam)(l-bud)(catalog)(arriba)(izq)(derecha)(cdn)(l-cdn)(l-exporta)] (nube) {};
  \node[anchor=north west, inner sep=0pt] at (nube.north west) {\includegraphics[width=6mm]{iconos/aws_cloud.pdf}};
  \node[anchor=north west, font=\scriptsize\bfseries] at ($(nube.north west)+(6.5mm,-0.5mm)$) {Nube de AWS (us-east-1)};
\end{scope}
\end{tikzpicture}
}
\caption{Servicios de AWS que usa el sistema (íconos oficiales~\cite{aws_icons}). La línea continua sigue a los datos y la punteada, a las órdenes. Cada lunes, Step Functions corre la cadena y termina exportando los resúmenes de oro a la web, que CloudFront entrega a la ciudadanía.}
\label{cloud:fig:arquitectura}
\end{figure*}

\begin{table}[!htbp]
\caption{Qué guarda cada capa y en qué formato}\label{cloud:tab:capas}
\centering
\footnotesize
\begin{tabular}{@{}Q{0.13\columnwidth}Q{0.36\columnwidth}Q{0.40\columnwidth}@{}}
\toprule
\textbf{Capa} & \textbf{Formato} & \textbf{Contenido} \\
\midrule
Bronce & JSON Lines + gzip, por carga & Lo que entrega la API, tal cual \\
Plata  & Iceberg, por mes de firma & 16 columnas limpias y sin duplicados \\
Oro    & Iceberg & Seis tablas agregadas por objetivo \\
\bottomrule
\end{tabular}
\end{table}

\subsubsection{Decisiones de diseño}
La descarga va por meses, y por días en los meses más grandes, para que ninguna Lambda pase de 15 minutos~\cite{aws_lambda_quotas}. Un estado \texttt{Map} corre tres tramos a la vez~\cite{aws_sfn_map} y cada Lambda baja cuatro páginas en paralelo con \texttt{ThreadPoolExecutor}~\cite{psf_futures}. Las consultas se ordenan por \texttt{:id} para que la paginación no salte ni repita contratos~\cite{tyler_order}.

Plata y oro están en Iceberg~\cite{asf_iceberg_spec}: cada escritura se confirma completa o no se confirma, se pueden actualizar filas con \texttt{MERGE}, se pueden añadir columnas después, y Athena consulta esas tablas sin pasos intermedios. Plata guarda solo 16 columnas: Athena escanea menos y los datos personales no salen de bronce. Para gastar poco, la Lambda corre en arm64 con 2.048~MB y Glue usa el modo Flex, que cuesta 0,29~USD por DPU-hora frente a 0,44~USD del modo estándar~\cite{aws_glue_pricing}.

\subsubsection{Seguridad}
El bucket es privado, tiene bloqueado el acceso público y cifra los datos con SSE-S3~\cite{aws_s3_sse}. A cada servicio se le dio un rol de IAM con los permisos estrictamente necesarios~\cite{aws_iam_bp}; la Lambda, por ejemplo, solo puede escribir dentro de \texttt{bronze/} (Fig.~\ref{cloud:lst:iam}). El token de la API está en las variables de entorno de la Lambda y no en el repositorio.

\begin{figure}[!htbp]



\begin{lstlisting}[language=json]
{
  "Version": "2012-10-17",
  "Statement": [{
    "Effect": "Allow",
    "Action": "s3:PutObject",
    "Resource": "arn:aws:s3:::<BUCKET_LAKE>/bronze/*"
  }]
}
\end{lstlisting}
\caption{Política \texttt{escribir-bronce} del rol de la Lambda}\label{cloud:lst:iam}
\end{figure}

% -----------------------------------------------------------------------------
\subsection{Implementación}\label{cloud:sec:impl}
Los recursos se configuraron a mano en us-east-1, paso a paso según el \texttt{README} del repositorio. La Tabla~\ref{cloud:tab:recursos} lista los recursos y la Fig.~\ref{cloud:lst:sfn} muestra el estado \texttt{Map} que reparte la carga histórica.

\begin{figure}[!htbp]



\begin{lstlisting}[language=json]
"DescargarTramos": {
  "Type": "Map",
  "ItemsPath": "$.windows",
  "MaxConcurrency": 3,
  "ItemProcessor": { "StartAt": "InvocarLambda",
    "States": { "InvocarLambda": {
      "Type": "Task",
      "Resource": "arn:aws:states:::lambda:invoke",
      "Parameters": { "FunctionName": "secop-ingest",
                      "Payload.$": "$" },
      "Retry": [{ "ErrorEquals": ["States.ALL"],
        "IntervalSeconds": 30, "MaxAttempts": 3,
        "BackoffRate": 2 }],
      "End": true } } }
}
\end{lstlisting}
\caption{Fragmento de \texttt{solo\_bronce.asl.json}}\label{cloud:lst:sfn}
\end{figure}

\begin{table}[!htbp]
\caption{Lo que quedó creado en la cuenta de AWS}\label{cloud:tab:recursos}
\centering
\footnotesize
\begin{tabular}{@{}Q{0.28\columnwidth}Q{0.68\columnwidth}@{}}
\toprule
\textbf{Recurso} & \textbf{Configuración} \\
\midrule
Bucket S3 del lago & Privado y cifrado; prefijos \texttt{bronze/}, \texttt{silver/}, \texttt{gold/} \\
6 roles de IAM & Las dos Lambdas, Step Functions, Scheduler, Glue y Budgets, cada uno con lo mínimo \\
Lambda \texttt{secop-ingest} & Python 3.12, arm64, 2.048~MB, 900~s, 4 hilos, páginas de 10.000 \\
Lambda \texttt{secop-web-export} & Python 3.12, arm64, 512~MB; 8 consultas a oro en Athena \\
Step Functions & \path{secop-solo-bronce} (histórico) y \path{secop-pipeline-diario} (hoy semanal; termina exportando la web) \\
EventBridge Scheduler & Cada lunes a las 6:00 (hora de Colombia), 2 reintentos \\
Glue 5.0 (Spark 3.5) & Plata con 2 o 10 \textit{workers}, oro con 5; Flex, concurrencia 1~\cite{aws_glue_iceberg} \\
Athena, \emph{workgroup} \texttt{secop} & Máximo 1~GB escaneado por consulta \\
AWS Budgets & 15~USD/mes, alertas al 50/80/100~\% y bloqueo automático \\
Web (S3 y CloudFront) & Bucket privado que solo lee CloudFront (OAC), HTTPS y caché de 15~min \\
\bottomrule
\end{tabular}
\end{table}

% -----------------------------------------------------------------------------
\subsection{Costos}\label{cloud:sec:costos}
La Tabla~\ref{cloud:tab:costos} combina los precios públicos de us-east-1~\cite{aws_pricing} con lo que se midió en el proyecto.

\begin{table}[!htbp]
\caption{Costos (USD). La columna Mes supone una actualización por semana; cuando era diaria, el mes costaba unos 2,4~USD.}\label{cloud:tab:costos}
\centering
\footnotesize
\begin{tabular}{@{}Q{0.18\columnwidth}Q{0.50\columnwidth}rr@{}}
\toprule
\textbf{Servicio} & \textbf{Base de cálculo} & \textbf{Hist.} & \textbf{Mes} \\
\midrule
Lambda & 12.800 GB-s, dentro de la capa gratuita~\cite{aws_lambda_pricing} & 0,00 & 0,00 \\
Step Functions & 0,025 USD por 1.000 transiciones~\cite{aws_sfn_pricing} & 0,02 & 0,01 \\
Glue (Flex) & 0,29 USD/DPU-h; histórico 10.274 DPU-s; 886 DPU-s por actualización & 0,83 & 0,31 \\
S3 & 0,023 USD/GB-mes, unos 3 GB~\cite{aws_s3_pricing} & 0,05 & 0,10 \\
Athena & 5 USD/TB escaneado~\cite{aws_athena_pricing} & 0,05 & 0,10 \\
CloudWatch & Registros de los servicios & 0,00 & 0,05 \\
Web & S3, CloudFront y Lambda dentro de la capa gratuita & -- & 0,00 \\
\midrule
\textbf{Total} & & \textbf{$\approx$1,0} & \textbf{$\approx$0,6} \\
\bottomrule
\end{tabular}
\end{table}

Para que el gasto no se dispare hay varios frenos. Un presupuesto de 15~USD al mes envía correos al 50, 80 y 100~\%, y al llegar al 100~\% una acción automática~\cite{aws_budget_actions} aplica una política que impide iniciar Step Functions, invocar la Lambda y lanzar Glue, sin borrar ningún dato. Los trabajos de Glue tienen tiempo máximo, una sola ejecución a la vez y ningún reintento automático. Athena corta toda consulta que vaya a leer más de 1~GB; contar los contratos de plata leyó solo 3,2~MB. Y si la carga semanal llega vacía, Glue no arranca.

% -----------------------------------------------------------------------------
\subsection{Resultados}\label{cloud:sec:avance}
Todo está desplegado y funcionando, incluida la web de resultados (Sección~\ref{sec:web}). Bronce tiene 5.687.382 contratos en 732 archivos (2,3~GB) y no se perdió ningún tramo. Plata guarda 5.687.110 contratos únicos en 137 particiones mensuales (107~MiB); se construyó en 199~s y costó 0,12~USD. De las seis tablas de oro (Tabla~\ref{cloud:tab:oro}), la de alertas ocupa casi todo el espacio, y las otras cinco son resúmenes pequeños, que son los que publica la web. Cuando hubo que añadir tres columnas, Iceberg las agregó sin reescribir los datos~\cite{asf_iceberg_evolution}, y reprocesar todo costó 0,44~USD.

\begin{table}[!htbp]
\caption{Tamaño de las seis tablas de oro}\label{cloud:tab:oro}
\centering
\footnotesize
\begin{tabular}{@{}lcrr@{}}
\toprule
\textbf{Tabla} & \textbf{Obj.} & \textbf{Filas} & \textbf{Col.} \\
\midrule
\path{gold_panorama_contratacion}   & 1 & 116.145 & 9 \\
\path{gold_alertas_contratos}       & 2 & 5.687.110 & 20 \\
\path{gold_riesgo_entidades}        & 2 & 5.119 & 20 \\
\path{gold_riesgo_departamentos}    & 2 & 34 & 17 \\
\path{gold_ley_garantias_diario}    & 3 & 2.219 & 7 \\
\path{gold_ley_garantias_entidades} & 3 & 4.984 & 9 \\
\bottomrule
\end{tabular}
\end{table}

\subsubsection{Consultas en Athena}
El 83,7~\% de los contratos activa al menos una alerta (Tabla~\ref{adq:t7}), y mientras regía la restricción electoral hubo 4.690 contratos directos cuya excepción no aparece en los datos (Tabla~\ref{adq:t8}). Encabeza la lista la Escuela de Formación del INPEC, con 266 contratos, seguida por la Alcaldía de Santa Marta con 187, la Gobernación de San Andrés con 174, la Agencia Nacional de Tierras con 171 y el Municipio de Soledad con 165. Como la causal pudo quedar mal registrada, estos contratos se revisan antes de sacar conclusiones.

Hay contratos con valores de hasta 4,4$\times 10^{17}$~COP; un hospital de Lérida aparece con 100 billones en un solo contrato. La regla de valor los deja fuera de las sumas sin borrarlos. Todos los casos revisados por encima de 10 billones eran errores. Entre 1 y 10 billones puede haber alguno real, como uno de la Registraduría por 3,3 billones, que también queda fuera.

\subsubsection{Carga semanal}
Cada lunes, Step Functions baja los últimos 90 días, comprueba si llegaron contratos y, si hay alguno, corre plata y oro y exporta la web. Hasta el 23 de noviembre de 2026 la carga era diaria y repetía 10 días; en la prueba del 13 de noviembre tardó 8~min~11~s, sumó 15.160 contratos nuevos y costó unos 0,07~USD. Para reducir costos pasó a una vez por semana: unos 0,3~USD al mes en vez de 2,1.

\subsubsection{Problemas y soluciones}
Con páginas de 50.000 contratos la Lambda se quedaba sin memoria, así que se bajaron a 10.000 y la memoria se subió a 2.048~MB. La Lambda venía configurada con 3~s de tiempo máximo y hubo que llevarla a 900~s. Cinco meses tardaban más de 15 minutos y se partieron en días. Los cortes de conexión de la API se cubrieron con reintentos en la Lambda y en Step Functions. La consola sugería Glue 6.0 con 480 minutos de límite; se dejó Glue 5.0, porque el código está escrito para Spark 3.5~\cite{aws_glue_versions}, con un límite mucho menor. La API publica con días de atraso, y una verificación contra la API mostró que algunos contratos aparecen uno o dos meses después de firmados (faltaban 6.434, casi todos de septiembre de 2026). Por eso la carga semanal repite los últimos 90 días y el \texttt{MERGE} evita los duplicados.

% -----------------------------------------------------------------------------
\subsection{Conclusiones}\label{cloud:sec:conclusiones}
\begin{itemize}
  \item El histórico completo quedó en AWS sin perder tramos y con un freno automático de gasto.
  \item Cuando cambiaron los objetivos, Iceberg permitió añadir columnas y reprocesar todo por 0,44~USD.
  \item Glue costó 0,83~USD contando las pruebas, y la carga semanal funciona sola por unos 0,07~USD por ejecución.
  \item La web de resultados ya está publicada; falta borrar los \emph{snapshots} viejos de Iceberg.
\end{itemize}
